% =====================================================================
% Section 07, Phase 1 research questions Q1–Q12
% =====================================================================
\makesectioncover{Phase 1 \textbullet\ The Twelve Questions}

\section{Phase 1 \textbullet\ The Twelve Research Questions}
\label{sec:phase1-questions}

The analysis notebook (\fn{Visual Analysis.ipynb}) answers twelve
questions in sequence. We chose each one to test a different
facet of the formal-vs-informal hypothesis or it isolates a single
structural gap. We document each question with the same five-part
template: \emph{Why}, \emph{Data}, \emph{Method}, \emph{Visualisation},
\emph{Finding}.

% --------------------------------------------------------------------
\subsection{Q1 \textbullet\ Job accessibility \texttimes\ alighting}

\textbf{Why.} If Cairo's transport network were demand-aligned, stops
with the highest daily alighting should also be in hexes with the
highest job accessibility, people get off where the work is. We test
that.

\textbf{Data.} \fn{merged\_A\_boarding\_population.csv} (every stop
attached to its containing population hex). 1{,}302 rows.

\textbf{Method.} Pearson and Spearman correlations between
\fn{Total\_Daily\_Alighting} and \fn{Jobs\_Accessible\_60min}; a
secondary 2-D density-contour visualisation to expose the joint
distribution rather than just a coefficient.

\figitem{phase1/Q1a_Where_Do_Cairo_s_Top_Alighting_Stops_Fall_on_Job_Accessibility.png}{1.0}{Q1a, Alighting vs job accessibility, top stops highlighted.}{fig:q1a}

\textbf{Why this viz.} A scatter shows the relationship strength,
dispersion, and outliers all at once; a 2-D density contour
(Q1b) backs it up by showing where the bulk of stops actually live in
the joint distribution rather than just the ones that sit far from
the mean.

\textbf{Finding.} \stat{$r = 0.261$} morning, $r = 0.248$ daily, 
\emph{weak positive} correlation. There is some demand-following but
not enough to claim Cairo's stop placement is built around the job
geography. This is the first quantitative crack in the
formal-network-as-demand-aligned hypothesis.

\figitem{phase1/Q1b_2D_Density_Job_Accessibility_Daily_Alighting.png}{0.95}{Q1b, Joint density of job accessibility and daily alighting.}{fig:q1b}

% --------------------------------------------------------------------
\subsection{Q2 \textbullet\ Terminal route symmetry}

\textbf{Why.} A balanced terminal is one where roughly the same number
of routes start there as end there. Strong asymmetry suggests
operational dysfunction (one-way runs, scheduling gaps).

\textbf{Data.} \fn{merged\_E\_terminal\_pop\_flow.csv}, terminals with
flow context.

\textbf{Method.} For each terminal compute
$|R_\text{origin} - R_\text{dest}| / R_\text{total}$; rank descending;
visualise as paired horizontal bars and as an asymmetry distribution.

\figitem{phase1/Q2a_Terminal_Symmetry_Origin_vs_Destination_Routes.png}{1.0}{Q2a, Per-terminal origin vs destination route counts.}{fig:q2a}

\textbf{Why this viz.} Horizontal paired bars make the
origin-vs-destination ratio readable at a glance for hundreds of
terminals; the second chart (Q2b) summarises the population-level
distribution of asymmetry.

\textbf{Finding.} A long right tail of asymmetric terminals (top 20
have more than 80\% of their routes on a single side). Asymmetry is
\emph{not} concentrated at low-traffic terminals, several major hubs
also exhibit it. This is one of the early signs that the network's
operational design is not purely demand-driven.

\figitem{phase1/Q2b_How_Common_Is_Asymmetry_and_Which_Terminals_Show_It_Most.png}{0.95}{Q2b, Distribution of route asymmetry across all 280 terminals.}{fig:q2b}

% --------------------------------------------------------------------
\subsection{Q3 \textbullet\ Drop points within / beyond the route corridor}

\textbf{Why.} If route lines are honest, almost all alighting should
happen \emph{within a short walking distance of the line itself}. This
is a basic adherence check, does the published route describe the
real-world corridor?

\textbf{Data.} \fn{q3\_drop\_points.csv}, boarding stops flagged
\fn{Near\_Route} when within 120\,m of the union of all route lines.

\textbf{Method.} Buffer the union of all route geometries by 120\,m;
flag every boarding point as inside or outside; aggregate.

\figitem{phase1/Q3a_Alighting_Near_vs_Far_from_Route_Corridors.png}{1.0}{Q3a, Alighting near vs. far from designated route corridors.}{fig:q3a}

\textbf{Why this viz.} Two clear bars (near vs far) is the right form
when the question is binary; a histogram of alighting volume by
proximity (Q3b) shows whether the rule holds across the full demand
distribution.

\textbf{Finding.} \stat{$88.0\%$ of alighting} happens within 120\,m of
the designated corridor. The published corridor is a faithful
description of where people actually get off. This is one of the few
\emph{validating} results in Phase 1, it gives us confidence to use
the route geometry as a meaningful spatial reference for downstream
buffer analyses (Q11, BRT corridor in Phase 2).

\figitem{phase1/Q3b_Distribution_of_Alighting_by_Route_Proximity.png}{0.95}{Q3b, Distribution of alighting volume by proximity bucket.}{fig:q3b}

% --------------------------------------------------------------------
\subsection{Q4 \textbullet\ Formal vs informal mode dominance}

\textbf{Why.} The defining structural fact of Cairo transport is that
informal microbuses carry the modal share. Q4 quantifies that share
and shows where it shifts (or doesn't) between morning peak and the
full day.

\textbf{Data.} \fn{q4\_formal\_vs\_informal.csv}.

\textbf{Method.} Per-stop ratios
\fn{Formal\_Percentage\_Morning} and
\fn{Formal\_Percentage\_Daily}; per-stop dominance flags;
\fn{Dominance\_Shifted} when morning-dominance differs from
daily-dominance.

\figitem{phase1/Q4a_Distribution_of_Formal_Transport_Percentage_Morning.png}{0.95}{Q4a, Distribution of formal-mode percentage at morning peak.}{fig:q4a}
\figitem{phase1/Q4b_Morning_vs_Daily_Formal_percent_does_preference_shift.png}{0.95}{Q4b, Morning vs daily formal percentage. Above the diagonal: stops where formal share grows during the day. Below: stops where formal share shrinks.}{fig:q4b}
\figitem{phase1/Q4c_Total_Boardings_by_Mode_Time_of_Day_Formal_vs_Informal.png}{0.95}{Q4c, Total boardings by mode and time-of-day.}{fig:q4c}

\textbf{Why these visuals.} Q4 needs five charts because the question
has five flavours: distribution of mode share, time-of-day shift,
absolute volume, fare economics, and per-stop ranking. Each plot
isolates one. The diagonal scatter (Q4b) is the most analytically
informative, it lets the reader see at once whether
formal-vs-informal is fixed or shifts during the day.

\textbf{Finding.} Most stops sit \emph{below} the diagonal (formal
share \emph{declines} from morning to all-day), informal microbuses
are the dominant non-peak mode. Cairo's network does not have a
``commuter rush'' formal layer that recedes; informal carries
disproportionately at off-peak.

\figitem{phase1/Q4d_Fare_vs_Route_Length_Formal_vs_Informal_Economics.png}{0.95}{Q4d, Fare vs route length, faceted by transport class. Slope = LE/km.}{fig:q4d}
\figitem{phase1/Q4e_Top_20_Stops_by_Informal_Demand_Formal_vs_Informal_Side-by-Side.png}{0.95}{Q4e, Top-20 stops by informal demand, with their formal counterpart on the same axis.}{fig:q4e}

% --------------------------------------------------------------------
\subsection{Q5 \textbullet\ Population density vs terminals}

\textbf{Why.} A direct test of whether the existing terminal layer
follows population.

\textbf{Data.} \fn{q5\_hex\_terminal\_counts.csv}, every population
hex with a count of terminals inside.

\textbf{Method.} Spearman correlation; box plot of population in hexes
with vs without a terminal.

\figitem{phase1/Q5a_Population_vs_Terminal_Count_per_Hexagon.png}{1.0}{Q5a, Population per hex vs number of terminals.}{fig:q5a}
\figitem{phase1/Q5b_Population_Distribution_Hexagons_With_vs_Without_Terminals.png}{0.9}{Q5b, Population distribution: hexes with terminals vs without.}{fig:q5b}

\textbf{Why this viz.} A scatter exposes outliers (terminals in
empty hexes); the box plot complements it by reducing the same data
to a clean median comparison.

\textbf{Finding.} Hexes \emph{with} a terminal have a higher median
population than those without, but the overlap is large. Terminals
are loosely demand-aligned in aggregate, but the per-cell mismatch is
the gap that Q9 isolates.

% --------------------------------------------------------------------
\subsection{Q6 \textbullet\ Route supply vs passenger demand}

\textbf{Why.} Are stops with many routes the stops that carry many
passengers? In a healthy network this should be a clean positive
relationship.

\textbf{Data.} \fn{q6\_routes\_vs\_passengers.csv}.

\textbf{Method.} Spearman correlation, an OLS trend line, and a
box-plot binning routes-per-stop into supply tiers.

\figitem{phase1/Q6a_Routes_Supply_vs_Passenger_Demand.png}{1.0}{Q6a, Per-stop route count vs passenger activity. The trendline shows the expected positive slope.}{fig:q6a}
\figitem{phase1/Q6b_Average_Passenger_Activity_by_Route_Count_Bin.png}{0.95}{Q6b, Average passenger activity per route-supply tier.}{fig:q6b}
\figitem{phase1/Q6c_Major_High-Demand_Stops.png}{0.95}{Q6c, The high-demand stops singled out for closer inspection.}{fig:q6c}

\textbf{Why this viz.} A scatter with trendline answers ``does the
relationship exist''; the binned box plot answers ``how strong / how
clean is it''.

\textbf{Finding.} Positive but noisy. The notebook surfaces several
high-demand stops with unexpectedly few routes (and vice-versa), 
these are the ``overcrowded / underserved'' terminals later
re-classified in Q11.

% --------------------------------------------------------------------
\subsection{Q7 \textbullet\ Morning vs evening symmetry}

\textbf{Why.} A residential commuter pattern is asymmetric (people
leave residential areas in the morning, come back in the evening). A
commercial hub is the mirror. Q7 separates the two.

\textbf{Data.} \fn{q7\_morning\_vs\_evening.csv}.

\textbf{Method.} A log-log plot of morning vs evening boarding makes
both regimes visible at once; a histogram of the imbalance ratio
quantifies how skewed the population is.

\figitem{phase1/Q7a_Log-Scale_Transit_Flow_Residential_vs_Commercial_Hubs.png}{1.0}{Q7a, Log-log of morning vs evening boarding. Above the diagonal: residential. Below: commercial.}{fig:q7a}
\figitem{phase1/Q7b_Morning_Evening_Imbalance_Distribution.png}{0.95}{Q7b, Distribution of the morning/evening imbalance ratio.}{fig:q7b}

\textbf{Why this viz.} Log-log compresses the dynamic range of stop
volumes (which span 4 orders of magnitude) and the diagonal becomes a
visually clean reference for ``balanced''.

\textbf{Finding.} The Cairo stop population is bimodal, clear
clusters of residential stops above the diagonal and commercial stops
below. The imbalance histogram shows long tails on both sides, most
stops are clearly one type or the other.

% --------------------------------------------------------------------
\subsection{Q8 \textbullet\ Spatial autocorrelation (Moran-style)}

\textbf{Why.} Are high-boarding stops randomly scattered or
geographically clustered? If they cluster, the network has
``regimes'', neighbourhoods of similar transport behaviour, and
those regimes are exactly what an under-served-area metric should
identify.

\textbf{Data.} \fn{q8\_spatial\_autocorrelation.csv}, each boarding
stop with its 8-nearest-neighbours mean boarding.

\textbf{Method.} Compute neighbour-mean boarding via KNN ($k=9$,
self-included then dropped); plot value vs neighbour-mean on log-log
axes; quadrant-classify into HH (high-high cluster), HL (high outlier),
LH (low outlier), LL (low cluster).

\figitem{phase1/Q8a_Enhanced_Moran_Scatter_Detecting_Spatial_Transit_Regimes.png}{1.0}{Q8a, Moran-style local-lag scatter. The HH and LL regions describe spatial clusters of similar stops.}{fig:q8a}
\figitem{phase1/Q8b_Network_Composition_Distribution_of_Spatial_Clusters.png}{0.85}{Q8b, Composition of the 4 spatial regimes.}{fig:q8b}

\textbf{Why this viz.} A Moran-style scatter is the canonical visual
test for spatial autocorrelation; the quadrant overlay does the
classification work in-figure rather than in a separate table.

\textbf{Finding.} Clear positive autocorrelation, HH and LL clusters
dominate. Cairo's transport behaviour comes in neighbourhoods, not in
randomly-distributed stops. (We will revisit this exact phenomenon in
Phase 2's H1 with a formal Moran's~I p-value.)

% --------------------------------------------------------------------
\subsection{Q9 \textbullet\ Underserved hexagons (Gap G4)}

\textbf{Why.} If the network were demand-following, every dense
population hex should see a proportional amount of boarding. Q9
defines a per-hex score
\fn{Underserved\_Score = pop\_rank(pct) - boarding\_rank(pct)} and
ranks the worst.

\textbf{Data.} \fn{q9\_underserved.csv}.

\textbf{Method.} Percentile-rank both population and boarding within
each hex; subtract; clip to non-negative; rank descending.

\figitem{phase1/Q9a_Priority_Analysis_High-Impact_Transit_Deserts.png}{1.0}{Q9a, High-impact transit deserts: hexes with a large positive Underserved Score.}{fig:q9a}

\textbf{Why this viz.} A two-axis priority scatter (population on x,
underserved score on y) lets the reader see size and severity in one
figure; the top-right quadrant is the policy target.

\textbf{Finding.} \stat{79 hexes} cross the \fn{Underserved\_Score $>$ 0.5}
threshold. They cluster in Imbaba, Shubra, Matariyya, the same
coral cluster the Phase 2 hero map will highlight. \kw{This is Gap G4.}

% --------------------------------------------------------------------
\subsection{Q10 \textbullet\ Empty-return terminals (Gap G2)}

\textbf{Why.} Vehicles that arrive at a terminal but leave empty are
the most basic operational waste signal. Q10 quantifies it.

\textbf{Data.} \fn{merged\_D\_terminal\_flow\_ratios.csv}.

\textbf{Method.} \fn{Empty\_Return\_Index = 1, passengers/vehicles},
clipped to $[0, 1]$. Severity score = index $\times \log(1 + \text{vehicles})$ to
combine rate \emph{and} volume.

\figitem{phase1/Q10_The_Empty_Return_Index_is_Binary_Healthy_vs_Completely_Wasteful.png}{1.0}{Q10, The Empty Return Index is binary in practice: 208 healthy vs 19 completely empty.}{fig:q10}
\figitem{phase1/Q10a_Top_20_Terminals_by_Empty_Return_Severity.png}{0.95}{Q10a, Top-20 terminals by empty-return severity.}{fig:q10a}
\figitem{phase1/Q10b_Vehicles_vs_Passengers_--_below_the_diagonal_empty_trips.png}{0.95}{Q10b, Vehicles vs passengers per terminal. Points below the diagonal are running below capacity.}{fig:q10b}

\textbf{Why this viz.} The bimodal histogram (Q10) is the single most
informative visual in Phase 1: it shows that the problem is \emph{not
gradual}. Terminals are either healthy (index = 0.0, 208 of them) or
completely broken (index $\geq$ 0.6, 19 of them), there is 
nothing in between.

\textbf{Finding.} \stat{19 critical empty-return terminals}, running
193 empty vehicle-trips every hour. Zero passengers. \kw{This is Gap G2.}

% --------------------------------------------------------------------
\subsection{Q11 \textbullet\ Ghost terminals (Gap G1) and zone classification}

\textbf{Why.} The single most important Phase 1 question. A ``ghost''
terminal has at least one route but \emph{zero boarding} within a 100\,m
walking radius. Q11 identifies them, classifies them by recovery
distance, and zones every terminal in the city.

\textbf{Data.} \fn{merged\_G\_underused\_terminals.csv}.

\textbf{Method.} 100\,m boarding buffer per terminal; flag underused;
animate the buffer radius from 100\,m to 1\,km to show how many
recover; treemap the route share of each operational zone.

\figitem{phase1/Q11a_Terminal_Utilization_at_100m_Baseline.png}{0.95}{Q11a, Terminal utilisation at 100 m baseline. The 115 ghosts are the red points.}{fig:q11a}
\figitem{phase1/Q11_Interactive_Slider_Walking_Radius_--_Watch_Red_Ghost_Terminals_Turn_Green.png}{0.95}{Q11 (slider), Walking radius from 100\,m to 1\,km. The interactive notebook version animates this; here we show the 500\,m frame.}{fig:q11slider}
\figitem{phase1/Q11_Recovery_Curve_--_How_Many_Ghost_Terminals_Disappear_as_Walking_Radius_Grows.png}{0.95}{Q11, Recovery curve: how many ghosts disappear as the walking radius grows.}{fig:q11recovery}
\figitem{phase1/Q11_What_Kind_of_Ghost_Terminal_--_69_need_stop_relocation_only_29_need_route_cancellation.png}{0.95}{Q11, Ghost-terminal classification: 69 near-miss (relocate) · 17 stop-too-far · 29 truly isolated.}{fig:q11kind}
\figitem{phase1/Q11b_Operational_Zone_Map_--_Strategy_Overview_at_100m.png}{0.95}{Q11b, Operational zones (Efficient Hub, Marginal, Overcrowded/Underserved, Dead Weight).}{fig:q11b}
\figitem{phase1/Q11c_Zone_Comparison_--_Same_Route_Budget_Vastly_Different_Output.png}{0.95}{Q11c, Same route budget, vastly different boarding output. The Dead Weight bar is the wasted-infrastructure signature.}{fig:q11c}
\figitem{phase1/Q11d_Route_Budget_Treemap_--_Box_area_share_of_total_routes.png}{0.95}{Q11d, Route-budget treemap. Box area = share of all 1{,}784 routes attached to that zone.}{fig:q11d}

\textbf{Why these visuals.} The seven Q11 charts together convert one
spatial question into a complete operational diagnostic. The slider /
recovery curve answers ``how many of these are real ghosts vs
walkability problems''; the zone treemap answers ``how much of the
network's route budget is actually living in dead-weight terminals''.

\textbf{Finding.} \stat{115 ghost terminals} at the 100\,m baseline,
broken into:
\begin{itemize}
 \item \stat{69} near-miss (100--500\,m), relocate the stop;
 \item \stat{17} stop-too-far (500\,m--1\,km), walkability fix;
 \item \stat{29} truly isolated ($>$1\,km), audit / decommission.
\end{itemize}
The recovery curve shows that 60\% disappear once the radius grows to
500\,m. The zone treemap reveals that ``Dead Weight'' terminals
consume the same route budget as Efficient Hubs but produce
\stat{1{,}870$\times$ fewer passengers}. \kw{This is Gap G1.}

% --------------------------------------------------------------------
\subsection{Q12 \textbullet\ Vehicle-route fit (Gap G3)}

\textbf{Why.} A microbus has 14 seats. Putting it on a 60-km route is
operationally absurd. Q12 quantifies how often it happens.

\textbf{Data.} \fn{q12\_vehicle\_type\_route\_length.csv} +
\fn{cleaned\_routes.csv}.

\textbf{Method.} Route length distribution by vehicle type (violin);
average length crossed with seat capacity (heat map); fare-per-km by
type (box); ECDF for the full distribution.

\figitem{phase1/Q12a_Route_Length_Distribution_by_Vehicle_Type_violin.png}{1.0}{Q12a, Route length distribution by vehicle type. Microbus carries the long-route tail.}{fig:q12a}
\figitem{phase1/Q12b_Avg_Route_Length_Vehicle_Type_x_Capacity.png}{0.95}{Q12b, Average route length per (vehicle type $\times$ capacity) cell. The mismatch is in the upper-left.}{fig:q12b}
\figitem{phase1/Q12c_Fare_per_km_by_Vehicle_Type_which_type_offers_best_value.png}{0.95}{Q12c, Fare per km by vehicle type. Microbus charges \stat{0.62 LE/km}, Tomnaya \stat{1.30}; bus and metro hover near \stat{0.22--0.25}.}{fig:q12c}
\figitem{phase1/Q12d_Vehicle_Type_Summary_Length_vs_Fare_vs_Capacity_bubble_size_seats.png}{0.95}{Q12d, Length vs fare with bubble-size = seats. Visual proof that the smallest vehicles run the longest routes.}{fig:q12d}
\figitem{phase1/Q12e_ECDF_Cumulative_Route_Length_Distribution_by_Vehicle_Type.png}{0.95}{Q12e, ECDF of route length by vehicle type.}{fig:q12e}

\textbf{Why these visuals.} Five charts, five facets: distribution
(violin), pivoted average (heat map), per-km cost (box), summary scatter
(bubble), cumulative form (ECDF). The bubble plot (Q12d) is the most
direct argument for the gap: bubble size = seats, so the smallest
bubbles consistently sit at the highest route-length values.

\textbf{Finding.} \stat{75\% of routes longer than 50 km use the
smallest vehicle (microbus)}; 44 specific long-route violations are
identified. The fare/km box puts numbers on it: informal modes charge
\stat{2.5--6$\times$ the per-km rate of the metro}. Riders pay the
informal premium because the formal alternative does not reach them.
\kw{This is Gap G3.}
